\documentstyle[%


righttag%


,11pt%


,amssymb%


,russian%               remove in Eng.


,izvru%                 Set izven% in Eng.


]{amsart}





% {\to}


\newcommand{\tg}{\operatorname{tg}}


\newcommand{\dint}{\displaystyle\int}


\newcommand{\tts}[1]{}





\theoremstyle{plain}


\theorembodyfont{\it}


\newtheorem{thmnonum}{\hskip\parindent Теорема}


\newtheorem{assrt}{\hskip\parindent Утверждение}


\newtheorem{hlt}{\hskip\parindent Теорема Харди--Литтлвуда}


\renewcommand{\thethmnonum}{}


\renewcommand{\thehlt}{}





%\hoffset=-15mm%                  For Eng.


\setcounter{page}{42}


\god{2003}


\nomer{5~(492)} %                        Remove in Eng.


%\nomer{Vol.\,47, No.\,5}%               Set in Eng.


%\str{\pageref{begin}--\pageref{end}}%  Set in Eng.


\udk{517.518}





\author{\it Б.В.\,Симонов}


% NB:   ^^ remove in Eng.


\title{О тригонометрических рядах


с $(K,S)$-монотонными коэффициентами


в пространствах с весом}


\thanks{Работа выполнена при финансовой поддержке Российского фонда


фундаментальных исследований, грант \No\,00-01-00042.}





\date{}


%\pagestyle{myheadings}%         Set in Eng.


\pagestyle{plain} %              Remove in Eng.


\begin{document}


%\thispagestyle{plain}%          Set in Eng.


\maketitle


\label{begin}











{\bf 1.} {\it Введение.} Известна





\begin{hlt}[{\series{m}\shape{n}\selectfont [1], c.\,657}]


Пусть $a_n\to 0$ при $n\to \infty$


и $a_n\geq a_{n+1}$ для всех $n$, где $\{a_n\}$ --- коэффициенты рядов


$\frac{a_0}{2}+\sum\limits_{n=1}^{\infty}a_n\cos nx$  и


$\sum\limits_{n=1}^{\infty}a_n\sin nx$,


а $f(x)$ и $g(x)$ --- суммы этих рядов.


Тогда для $p\in (1,\infty)$ справедливы неравенства


\begin{align*}


C_1\bigg(\sum\limits_{n=0}^{\infty}a_n^p(n+1)^{p-2}\bigg)^{\frac{1}{p}}


&\leq \|f(x)\|_p \leq C_2 \bigg(\sum\limits_{n=0}^{\infty}


a_n^p(n+1)^{p-2}\bigg)^{\frac{1}{p}},\\


%


C_3\bigg(\sum\limits_{n=1}^{\infty}a_n^pn^{p-2}\bigg)^{\frac{1}{p}}


&\leq \|g(x)\|_p \leq C_4 \bigg(\sum\limits_{n=1}^{\infty}


a_n^pn^{p-2}\bigg)^{\frac{1}{p}},


\end{align*}


где положительные постоянные $C_1,\dotsc,C_4$ не зависят от $\{a_n\}$.


\end{hlt}





Заметим, что все эти (и нижеследующие) неравенства 


понимаются таким образом: из конечности правой части следует


конечность их левой части.\\





Приведем утверждение, дополняющее теорему


Харди-Литтлвуда.








\begin{thmnonum}[{\series{m}\shape{n}\selectfont [2]}]


{\em а)} Пусть $a_n\to 0$ при $n\to \infty$


и $a_n-2a_{n+1}+a_{n+2}\geq 0$ для всех $n$.





Тогда для $p\in (0,\infty)$


$$


C_1\bigg(\sum\limits_{n=0}^{\infty}(a_n-a_{n+1})^p(n+1)^{2p-2}


\bigg)^{\frac{1}{p}}


\leq \|f(x)\|_p \leq C_2 \bigg(\sum\limits_{n=0}^{\infty}(a_n-


a_{n+1})^p(n+1)^{2p-2}\bigg)^{\frac{1}{p}}.


$$





{\em б)} Пусть $a_n\to 0$ при $n\to \infty$


и $a_n\geq a_{n+1}$ для всех $n$.





Тогда для $p\in (0,\infty)$


$$


C_3\bigg(\sum\limits_{n=1}^{\infty}a_n^pn^{p-2}\bigg)^{\frac{1}{p}}


\leq \|g(x)\|_p \leq C_4 \bigg(\sum\limits_{n=1}^{\infty}


a_n^pn^{p-2}\bigg)^{\frac{1}{p}},


$$


где положительные постоянные $C_1,\dots,C_4$ не зависят от $\{a_n\}$.


\end{thmnonum}





Нетрудно заметить, что при выполнениии условия а) теоремы


справедливы неравенства


$C_1(a_n-a_{n+2})^p\leq (a_n-a_{n+1})^p\leq C_2(a_n-a_{n+2})^p$,


где положительные постоянные $C_1$, $C_2$ 


не зависят от $n$ и $\{a_n\}$. 


На основании этого п.\,а) теоремы из [2] может быть записан в


эквивалентной форме


\begin{equation}


C_3\bigg(\sum\limits_{n=0}^{\infty}(a_n-a_{n+2})^p(n+1)^{2p-2}


\bigg)^{\frac{1}{p}}


\leq \|f(x)\|_p \leq 


C_4 \bigg(\sum\limits_{n=0}^{\infty}(a_n-


a_{n+2})^p(n+1)^{2p-2}\bigg)^{\frac{1}{p}}.


\end{equation}





Результаты приведенных выше теорем обобщает





\begin{assrt}


Теорема Харди--Литтлвуда и п.\,{\em б)} теоремы из $[2]$


имеют место и для тригонометрических


рядов, коэффициенты которых удовлетворяют свойствам


$a_n{\to} 0$ $(n{\to} \infty)$,


$a_{2n}=0$, $a_{2n+1}\geq a_{2n+3}$, $n=0,1,2,\dotsc;$


п.\,{\em а)} теоремы из $[2]$ имеет место и для тригонометрических рядов,


коэффициенты которых удовлетворяют свойствам


$a_n\to 0$ $(n\to \infty)$, $a_0=\frac{a_1-a_2}{2}$,


$a_{2n}=0$, $a_{2n-1}-2a_{2n+1}+ a_{2n+3}\geq 0$, $n=1,2,\dotsc$


\end{assrt}





Заметим, что условие монотонности 


и выпуклости для этих коэффициентов


не выполнено, за исключением случая, когда $a_{n} \equiv 0.$\\





Введем основные определения и обозначения для формулировки


результатов данной работы.


Пусть $\Phi$ --- совокупность измеримых, неотрицательных, суммируемых


на $(-\pi,\pi)$ функций. Пространством с весом $L((-\pi,\pi);p;


\varphi)$, где $p>0$, $\varphi \in \Phi$, назовем множество


измеримых, $2\pi$-периодических функций $f(x)$, для каждой


из которых конечна квазинорма


$$


\|f(x)\|_{((-\pi,\pi);p;\varphi)}=


\bigg(\int_{-\pi}^{\pi}\varphi(x)|f(x)|^p dx \bigg)^{\frac{1}{p}}.


$$





Будем рассматривать ряды


\begin{equation}


\frac{a_0}{2}+\sum\limits_{n=1}^{\infty} a_n\cos nx


\end{equation}


и


\begin{equation}


\frac{a_0}{2}+\sum\limits_{n=1}^{\infty} a_n\sin nx.


\end{equation}


Для целых $k\geq 0$ обозначим


\begin{alignat*}{2}


\Delta_k a_m&=\sum\limits_{i=0}^{k}(-1)^i C_k^i a_{m+i} &\quad 


(\Delta_0 a_m&=a_m), \\


%


\{ \Delta\}_k a_m &=\sum\limits_{i=0}^{k} C_k^i a_{m+i} &\quad 


(\{ \Delta\}_0a_m &=a_m).


\end{alignat*}


Пусть $\varphi\in\Phi$, $\varphi_c(x)=\varphi(x)+\varphi(-x)$.


Скажем, что функция $\varphi(x)$ удовлетворяeт 





{\it условию} $A_1$, если для любого $\delta\in (0,\pi)$ 


$$


\int_{0}^{\delta}\varphi_c(x) dx \leq


C_1\int_{\frac{\delta}{2}}^{\delta}\varphi_c(x) dx,


$$


где положительная постоянная $C_1$ не зависит от $\delta$;





{\it условию} $A_2$, если для любого $\delta\in (0,\pi)$


и любого $b\geq \frac{\pi}{2\delta}$\\


$$


\int_{\frac{\delta}{2}}^{\delta}\varphi_c(x)\sin^2 bx\, dx \geq


C_2\int_{\frac{\delta}{2}}^{\delta}\varphi_c(x) dx,


$$


где положительная постоянная $C_2$ не зависит от $b$ и $\delta$;





{\it условию} $A_3$, если $\varphi(-x)=\varphi(x-\pi)$,


$\varphi(x)=\varphi(\pi-x)$ для почти всех $x\in (0,\pi)$;





{\it условию} $A(p)$ (при заданном $p$), 


если для любого $\delta\in (0,\pi)$ 


$$


\int_{\frac{\delta}{4}}^{\pi}\varphi_c(x) x^{-p} dx \leq


C_3\int_{\frac{\delta}{2}}^{\delta}\varphi_c(x) x^{-p}dx,


$$


где положительная постоянная $C_3$ не зависит от $\delta$.





Приведем пример функции $\varphi (x)$, 


для которой условие $A_2$ заведомо выполнено.


Для этого введем дополнительно два определения.





Следуя С.Н.\,Бернштейну ([3]), назовем конечную функцию $\alpha(x)$


почти возрастающей на $(a,b)$, если существует такая постоянная


$C{>}0$, что $\alpha(x_1)\leq C\alpha(x_2)$ при любых $x_1$ и $x_2$ c


$a<x_1<x_2<b$.





Скажем, что конечная функция $\alpha (x)$ удовлетворяет


$\Delta_2$-условию на $(0,b)$, если существует такая постоянная


$C>0$, что


$\alpha (x) \leq C \alpha(x/2)$ для любого $x \in (0,b)$.


 


Теперь рассмотрим конечную функцию $\varphi (x)$, являющуюся четной


на $(-\pi,\pi)$, почти возрастающей на $(0,\pi)$ и удовлетворяющей


$\Delta_2$-условию на $(0,\pi)$. Для этой функции условие $A_2$


выполнено. Это легко получить, используя свойства функции $\varphi(x)$


и неравенство ([2])


\begin{equation}


\int_{a\pi}^{2a\pi} \frac{\sin^2x}{x} dx 


\geq C,


\end{equation}


где положительная постоянная $C$ не зависит от $a$ ($a\geq\frac{1}{4}$).





Скажем, что последовательность 


$\{a_n\}_{n=0}^{\infty}= \{a_n\}$ \ $(k,s)$-монотонна, если


$a_n \to 0$ при $n\to \infty$ и


$\Delta_k (\{ \Delta \}_s a_n) \geq 0$ для некоторых $k\geq 0$, $s\geq 0$


и всех $n$.





Нетрудно проверить, что если последовательность $\{a_n\}$


($a_n\to 0$ при $n\to \infty$) невозрастающая, то она $(1,s)$-монотонна


для любого $s=0,1,\dotsc$


Обратное не всегда верно. Например, рассмотрим последовательность


$\{a_n\}$ такую, что $a_n \to 0$ при $n\to \infty$ и


$a_{2n}=0$, $a_{2n+1}\geq a_{2n+3}$ при $n=0,1,\dotsc$


Тогда эта последовательность не является невозрастающей, но является


$(1,1)$-монотонной.





Ряд (2) или (3) называется рядом с $(k,s)$-монотонными коэффициентами,


если последовательность $\{a_n\}$, составленная из коэффициентов этих


рядов, является $(k,s)$-монотонной.





Через $[a]$ обозначим целую часть числа $a$. Пусть $f(x)$ --- сумма ряда (2),


$g(x)$ --- сумма ряда (3).





В данной работе находятся условия, при которых суммы тригонометрических


рядов (2) и (3)


с $(k,s)$-монотонными коэффициентами принадлежат 


пространствам с весом,


а также устанавливаются оценки квазинорм этих функций через


коэффициенты рядов (2) и (3).





Далее приведем результаты, которые будут обоснованы в последующих пунктах


данной статьи.


 


\begin{assrt}


Пусть $\varphi\in\Phi$ и удовлетворяет условию $A_1$, а


последовательность $\{a_n\}$ является $(k,s)$-монотонной.


Тогда при любом $p\in (0,\infty)$





{\em a)} если $s=0$, $k\geq 2$, то


\begin{equation}


\| f(x)\|_{((-\pi,\pi);p;\varphi)}\leq


C_1\bigg(\sum\limits_{n=1}^{\infty}(\Delta_1a_{n-1})^p n^{2p}


\int_{\frac{\pi}{n+1}}^{\frac{\pi}{n}} \varphi_c(x) dx\bigg)^{\frac{1}{p}};


\end{equation}





если $s=1$ или $s=2$, $k\geq 2$, то


\begin{multline}


\| f(x)\|_{((-\frac{\pi}{2},\frac{\pi}{2});p;\varphi)}\leq


C_2\bigg(\int_{\frac{\pi}{4}}^{\frac{\pi}{2}} \varphi_c(x) dx


(\Delta_1(\{\Delta\}_sa_0))^p+ \\ +


\sum\limits_{n=2s}^{\infty}


\int_{\frac{\pi}{n+1}}^{\frac{\pi}{n}} \varphi_c(x) dx\, n^{2p}


(\Delta_1(\{\Delta\}_sa_{n-s}))^p \bigg)^{\frac{1}{p}};


\end{multline}





{\em б)} если $s=0$, $k\geq 1$, то


\begin{equation}


\| F(x)\|_{((-\pi,\pi);p;\varphi)}\leq


C_3\bigg(\int_{\frac{\pi}{2}}^{\pi} \varphi_c(x) dx


|a_0|^p+\sum\limits_{n=1}^{\infty}


\int_{\frac{\pi}{n+1}}^{\frac{\pi}{n}} \varphi_c(x) dx\, n^p\cdot a_n^p


\bigg)^{\frac{1}{p}};


\end{equation}





если $s=1$ или $s=2$, $k\geq 1$, то


\begin{multline}


\| F(x)\|_{((-\frac{\pi}{2},\frac{\pi}{2});p;\varphi)}\leq


C_4\bigg(\int_{\frac{\pi}{4}}^{\frac{\pi}{2}} \varphi_c(x) dx


(\{\Delta\}_sa_0+|a_0|(s-1))^p+ \\ +


\sum\limits_{n=2}^{\infty}


\int_{\frac{\pi}{n+1}}^{\frac{\pi}{n}} \varphi_c(x) dx\, n^p


(\{\Delta\}_sa_{n-1})^p \bigg)^{\frac{1}{p}},


\end{multline}


где $F(x)=f(x)$ или $F(x)=g(x)$, а положительные постоянные


$C_1,\dotsc,C_4$ не зависят от последовательности $\{a_n\}$.


\end{assrt}





Отметим, что для справедливости неравенства (7)


в случае, когда $F(x)=g(x)$, достаточно потребовать $(k,s)$-монотонность


последовательности $\{a_n\}_{n=1}^{\infty}$.





\begin{assrt}


Пусть $\varphi\in\Phi$, а


последовательность $\{a_n\}$ является $(k,s)$-монотонной,


$p\in (0,\infty)$.





{\em а)} Если $s=0$, $k\geq 3$, то


\begin{equation}


\| f(x)\|_{((-\pi,\pi);p;\varphi)}\geq


C_1\bigg(\sum\limits_{n=1}^{\infty}(\Delta_1a_{n-1})^p n^{2p}


\int_{\frac{\pi}{n+1}}^{\frac{\pi}{n}} \varphi_c(x)dx\bigg)^{\frac{1}{p}};


\end{equation}





если $s=1$ или $s=2$, $k\geq 3$ $($в случае $s=2$ 


для ряда $(2)$ считаем дополнительно, что


$a_0\geq 5a_1+3a_2-8a_3-2a_4+3a_5)$, то


\begin{multline}


\| f(x)\|_{((-\frac{\pi}{2},\frac{\pi}{2});p;\varphi)}\geq


C_2\bigg(\int_{\frac{\pi}{4}}^{\frac{\pi}{2}} \varphi_c(x) dx


(\Delta_1(\{\Delta\}_sa_0))^p+\\ +


\sum\limits_{n=2s}^{\infty}


\int_{\frac{\pi}{n+1}}^{\frac{\pi}{n}} \varphi_c(x) dx\, n^{2p}


(\Delta_1(\{\Delta\}_sa_{n-s}))^p \bigg)^{\frac{1}{p}}.


\end{multline}


Кроме того, для функций $\varphi(x)$, дополнительно удовлетворяющих


условию $A(p)$, при $s=0$ и $k\geq 3$


\begin{equation}


\| f(x)\|_{((-\pi,\pi);p;\varphi)}\geq


C_3\bigg(\int_{\frac{\pi}{2}}^{\pi} \varphi_c(x) dx\,


a_0^p+\sum\limits_{n=1}^{\infty}


\int_{\frac{\pi}{n+1}}^{\frac{\pi}{n}} \varphi_c(x) dx\,n^p\cdot a_n^p


\bigg)^{\frac{1}{p}};


\end{equation}


при $s=1$ или $s=2$ и $k\geq 3$ 


\begin{equation}


\| f(x)\|_{((-\frac{\pi}{2},\frac{\pi}{2});p;\varphi)}\geq


C_4\bigg(\int_{\frac{\pi}{4}}^{\frac{\pi}{2}} \varphi_c(x) dx


(\{\Delta\}_sa_0)^p+ \sum\limits_{n=2}^{\infty}


\int_{\frac{\pi}{n+1}}^{\frac{\pi}{n}} \varphi_c(x) dx\,n^p


(\{\Delta\}_sa_{n-1})^p \bigg)^{\frac{1}{p}}.


\end{equation}





{\em б)} Eсли $s=0$, $k\geq 2$, то 


\begin{equation}


\| g(x)\|_{((-\pi,\pi);p;\varphi)}\geq


C_5\bigg(\int_{\frac{\pi}{2}}^{\pi} \varphi_c(x) dx\,


a_0^p+\sum\limits_{n=1}^{\infty}


\int_{\frac{\pi}{n+1}}^{\frac{\pi}{n}} \varphi_c(x) dx\,n^p\cdot a_n^p


\bigg)^{\frac{1}{p}};


\end{equation}


если $s=1$ или $s=2$, $k\geq 2$, то


\begin{equation}


\| g(x)\|_{((-\frac{\pi}{2},\frac{\pi}{2});p;\varphi)}\geq


C_6\bigg(\int_{\frac{\pi}{4}}^{\frac{\pi}{2}} \varphi_c(x) dx


(\{\Delta\}_sa_0)^p+\sum\limits_{n=2}^{\infty}


\int_{\frac{\pi}{n+1}}^{\frac{\pi}{n}} \varphi_c(x) dx\,n^p


(\{\Delta\}_sa_{n-1})^p \bigg)^{\frac{1}{p}};


\end{equation}


где положительные постоянные


$C_1,\dotsc,C_6$ не зависят от последовательности


$\{a_n\}$.


\end{assrt}





Отметим, что для справедливости неравенства (13)


достаточно потребовать $(k,s)$-мо\-но\-тон\-ность


последовательности $\{a_n\}_{n=1}^{\infty}$ и $a_0\geq 0$.





\begin{assrt}


Пусть $\varphi\in\Phi$ и удовлетворяет условию $A_2$, а


последовательность $\{a_n\}$ является $(k,s)$-монотонной.


Тогда при любом $p\in (0,\infty)$





{\em а)} если $s=0$, $k\geq 2$,


то справедливо неравенство $(9)$, при этом, если функция $\varphi(x)$


дополнительно удовлетворяет условию $A(p)$, 


то справедливо неравенство $(11);$





если $s=1$, $k\geq 2$, $a_0\geq a_2+a_1-a_3+4(a_2-a_4)$


для коэффициентов ряда $(2)$,


то справедливо неравенство $(10)$, кроме того, если функция $\varphi(x)$


дополнительно удовлетворяет условию $A(p)$, 


то справедливо неравенство $(12);$





если $s=2$, $k\geq 2$, $a_0\geq 3a_2+2a_1-2a_3-2a_4$


для коэффициентов ряда $(2)$,


то справедливо неравенство $(10)$, при этом, если функция $\varphi(x)$


дополнительно удовлетворяет условию $A(p)$, 


то справедливо неравенство $(12);$





{\em б)} если $s=0$, $k\geq 1$, то 


$$


\| g(x)\|_{((-\pi,\pi);p;\varphi)}\geq


C_1\bigg(\int_{\frac{\pi}{4}}^{\frac{\pi}{2}} \varphi_c(x) dx


(a_0^p+a_1^p)+\sum\limits_{n=2}^{\infty}


\int_{\frac{\pi}{n+1}}^{\frac{\pi}{n}} \varphi_c(x) dx\,n^p\cdot a_n^p


\bigg)^{\frac{1}{p}};


$$





если $s=0$, $k\geq 1$, $a_0\geq 0$


для коэффициентов ряда $(2)$, последовательность


$\{a_n\}_{n=1}^{\infty}$ \ $(k,s)$-монотонна, а функция $\varphi(x)$


дополнительно удовлетворяет условию $A_3$ и четная, то


справедливо неравенство $(13);$





если $s=1$, $k\geq 1$, то справедливо неравенство $(14);$





если $s=2$, $k\geq 1$, $a_0\geq 4a_1+a_2$


для коэффициентов ряда $(3)$, то


\begin{multline*}


\| g(x)\|_{((-\frac{\pi}{2},\frac{\pi}{2});p;\varphi)}\geq


C_2\bigg(\int_{\frac{\pi}{4}}^{\frac{\pi}{2}} \varphi_c(x) dx


((\{\Delta\}_2a_0)^p+(\{\Delta\}_2a_1)^p+ \\ +


(\{\Delta\}_2a_2)^p)+ 


\sum\limits_{n=4}^{\infty}


\int_{\frac{\pi}{n+1}}^{\frac{\pi}{n}} \varphi_c(x) dx\,n^p


(\{\Delta\}_2a_{n-1})^p \bigg)^{\frac{1}{p}},


\end{multline*}


где положительные постоянные


$C_1$, $C_2$ не зависят от последовательности $\{a_n\}$.


\end{assrt}





{\bf 2.} Пусть $B_0^1(x)=\frac{1}{2}$;


$B_n^1(x)=\frac{1}{2}+\cos x+\dotsb+\cos nx$ для $n\geq 1$;


$B_n^k(x)=\sum\limits_{\nu=0}^nB_{\nu}^{k-1}(x)$


для $k=2,3,\dotsc$ и $n\geq 0$;


$\overline{B}{}_n^1(x)=\sin x+\dotsb+\sin nx$ для $n\geq 1$;


$\overline{B}{}_n^k(x)=


\sum\limits_{\nu=1}^n\overline{B}{}_{\nu}^{k-1}(x)$ для


$k=2,3,\dotsc$ и $n\geq 1$.





\begin{lem}[{\series{m}\shape{n}\selectfont [4]}]


Пусть $a_n\geq 0$, $b_n\geq 0$, $n=1,2,\dotsc$,


$1\leq p<\infty$. Тогда





{\em a)} если $\sum\limits_{\nu=n}^{\infty} a_{\nu}= 


a_n\beta_n$, $n=1,2,\dotsc$, то 


$\sum\limits_{k=1}^{\infty}a_k\Big(\sum\limits_{\nu=1}^{k}b_\nu\Big)^p \leq


p^p\sum\limits_{\nu=1}^{\infty} a_\nu(\beta_\nu b_\nu)^p;$





{\em б)} если $\sum\limits_{\nu=1}^{n} a_{\nu}= 


a_n\beta_n$, $n=1,2,\dotsc$, то


$\sum\limits_{n=1}^{\infty}a_n\Big(


\sum\limits_{\nu=n}^{\infty}b_\nu\Big)^p \leq


p^p\sum\limits_{\nu=1}^{\infty}a_\nu(\beta_\nu b_\nu)^p$. 


\end{lem}





\begin{lem}[{\series{m}\shape{n}\selectfont [5], с.\,125}]


Пусть $a_n\geq 0$, $0<\alpha\leq \beta<\infty$. Тогда


$$


\bigg(\sum\limits_{n=1}^{\infty}a_n^\beta\bigg)^{\frac1\beta}


\leq \bigg(\sum\limits_{n=1}^{\infty}a_n^\alpha\bigg)^{\frac1\alpha}.


$$


\end{lem}





\begin{lem}[{\series{m}\shape{n}\selectfont [5], с.\,66}]


Пусть $u_1\geq 0$, $u_2\geq 0$. Тогда для любого


$\alpha >0$ \ $C_1(\alpha )(u_1^\alpha +u_2^\alpha )\leq


(u_1+u_2)^\alpha \leq C_2(\alpha )(u_1^\alpha +u_2^\alpha )$.


\end{lem}





\begin{lem}


Пусть последовательность $\{a_n\}$


является $(k,s)$-монотононной. Тогда





{\em а)} функция $f(x)$ может быть почти всюду представлена в виде


$$


f(x)=\frac{1}{2\sin\frac{x}{2}}\sum\limits_{n=1}^{\infty}\Delta _1 


a_{n-1}\sin(2n-1)\frac x2,


$$


если $k=1$, $s=0$.





если $k=1$, $s=1$, то


$$


f(x)=\frac{1}{2\cos\frac{x}{2}}\sum\limits_{n=1}^{\infty}\{\Delta \}_1 


a_{n-1}\cos(2n-1)\frac x2;


$$





если $k=1$, $s=2$, то


$$


f(x)=\frac{a_0-a_2}{2}\cdot\frac{1}{(2\cos\frac{x}{2})^2}+


\frac{1}{(2\cos\frac{x}{2})^2}\bigg(


\frac{\{\Delta \}_2 a_0}{2}+\sum\limits_{n=1}^{\infty}\{\Delta \}_2


a_{n-1}\cos nx \bigg);


$$





если $k=2$, $s=1$, то


\begin{multline*}


f(x)=\frac{a_0-a_2}{2}-


\frac{\Delta _1(\{\Delta \}_1 a_1)+4\Delta _1(\{\Delta \}_1


a_2)}{2(2\cos \frac{x}{2})^2}+


2\Delta _1(\{\Delta \}_1a_2)\tg^2\frac{x}{2}+\\


+\frac{1}{2\sin\frac{x}{2}(2\cos\frac{x}{2})^2}


((\Delta _1(\{\Delta\}_1 a_1)+\Delta _1(\{\Delta\}_1 


a_2))\bigg(\sin\frac x2+\sin 3\frac x2\bigg)+ \\ 


+\sum\limits_{n=3}^{\infty}\Delta _1(\{\Delta\}_1 


a_{n-2})\sin(2n-1)\frac{x}{2}\bigg);


\end{multline*}





если $k=2$, $s=2$, то


\begin{multline*}


f(x)=\frac{\Delta _2(\{\Delta\}_2a_0)}{2}-\frac{a_1-a_3}


{2(2\cos\frac{x}{2})^2}+


\frac{1}{(2\cos\frac{x}{2})^2}\bigg(


\frac{3a_1+4a_2-a_3-2a_4}{2} +\\


+(2a_1+3a_2-a_4)\cos x+


\sum\limits_{n=2}^{\infty}\{\Delta\}_2a_{n-1}\cos nx\bigg);


\end{multline*}





если $k=3$, $s=2$, то


\begin{multline*}


f(x)=\frac{\Delta _3(\{\Delta\}_2a_0)}{2}-\frac{2a_1-


3a_3+a_5}{2(2\cos\frac{x}{2})^2}+


\frac{1}{(2\cos\frac{x}{2})^2}\bigg(


\frac{6a_1+4a_2-7a_3-2a_4+3a_5}{2} + \\ +(3a_1+3a_2-2a_3-a_4+a_5)\cos x+


\sum\limits_{n=2}^{\infty}\{\Delta\}_2a_{n-1}\cos nx\bigg);


\end{multline*}





{\em б)} функция $g(x)$ может быть почти всюду представлена в виде





\begin{equation}


g(x)=\frac{a_0}{2}\bigg(1-\tg{\frac{x}{2}}\bigg)+


\frac{1}{(2\cos{\frac{x}{2}})^s}


\sum\limits_{n=1}^{\infty}\{\Delta\}_s a_{n-1}\sin (ns-2+s)


{\frac{x}{2}},


\end{equation}


если $k=1$, $s=1$ или $s=2$.


\end{lem}





{\bf Доказательство} проведем для п.\,б).





Случай $s=1$. Рассмотрим


$S_n(x)=\frac{a_0}{2}+\sum\limits_{j=1}^{n}a_j\sin jx$.


Тогда


\begin{gather*}


S_n(x)=\frac{a_0}{2}+\sum\limits_{j=1}^{n} a_j\sin 2j{\frac{x}{2}} 


\frac{\cos{\frac{x}{2}}}{\cos{\frac{x}{2}}}=


\frac{a_0}{2}+\sum\limits_{j=1}^{n} a_j\bigg(\sin (2j-1){\frac{x}{2}} +\\


%


+\sin(2j+1){\frac{x}{2}}\bigg)\frac{1}{2\cos{\frac{x}{2}}}=


\frac{a_0}{2}+\frac{1}{2\cos{\frac{x}{2}}}\bigg(a_1\sin{\frac{x}{2}}+


\sum\limits_{j=2}^{n}(a_{j-1}+a_j)


\sin(2j-1){\frac{x}{2}}\bigg)+\\


%


+\frac{1}{2\cos{\frac{x}{2}}}a_n\sin(2n+1){\frac{x}{2}}=


\frac{a_0}{2}\bigg(1-\tg{\frac{x}{2}}\bigg)+\\


%


+\frac{1}{\cos{\frac{x}{2}}}


\sum\limits_{j=1}^{n}(a_{j-1}+a_j)\sin(2j-1){\frac{x}{2}}


+\frac{1}{2\cos{\frac{x}{2}}}a_n\sin(2n+1){\frac{x}{2}}.


\end{gather*}


При $n\to \infty$ получаем представление функции $g(x)$.


Существование предела устанавливается так же, как в ([1], с.\,100).





Случай $s=2$.


Используя представление $S_n(x)$, полученное выше, будем иметь


\begin{gather*}


S_n(x)=\frac{a_0}{2}+ \frac{1}{2\cos{\frac{x}{2}}}


\bigg(a_1\sin{\frac{x}{2}}+


\sum\limits_{j=2}^{n}(a_{j-1}+a_j)\sin(2j-1){\frac{x}{2}}\bigg) 


\frac{\cos{\frac{x}{2}}}{\cos{\frac{x}{2}}}+\\


%


+\frac{1}{2\cos{\frac{x}{2}}}


a_n\sin(2n+1){\frac{x}{2}}=


\frac{a_0}{2}+\frac{1}{(2\cos{\frac{x}{2}})^2}\bigg((2a_1+a_2)\sin x+\\


%


+\sum\limits_{j=2}^{n}\{\Delta\}_2a_{j-1}\sin jx\bigg)


-(a_n+a_{n+1})\frac{\sin nx}{(2\cos{\frac{x}{2}})^2}+


\frac{1}{2\cos{\frac{x}{2}}} a_n \sin(2n+1){\frac{x}{2}} = \\


%


=\frac{a_0}{2}\bigg(1-\tg\frac{x}{2}\bigg)+\frac{1}{(2\cos\frac{x}{2})^2}


\sum\limits_{j=1}^{n}\{\Delta\}_2a_{j-1}\sin jx +


\frac{1}{2\cos{\frac{x}{2}}} a_n\sin(2n+1){\frac{x}{2}}- \\


%


-\frac{1}{(2\cos\frac{x}{2})^2}(a_n+a_{n+1})\sin nx.


\end{gather*}


При $n\to \infty$ получим представление функции $g(x)$ в этом случае.





Аналогичным образом доказывается п.\,а).$\quad\square$





\begin{lem}


Пусть $\varphi \in \Phi$ и удовлетворяет условию $A_1$,


а последовательность $\{b_n\}$ является $(1,0)$-монотонной. Тогда для


$p\in(0,\infty)$


\begin{gather}


\bigg(\int_{-\pi}^{\pi} \varphi(x) \bigg| \sum\limits_{n=1}^{\infty}


b_n \frac{\sin(2n-1)\frac{x}{2} }{\sin{\frac{x}{2}}}\bigg|^p dx 


\bigg)^{\frac{1}{p}}\leq 


C_1\bigg( \sum\limits_{n=1}^{\infty} b_n^p n^{2p}


\int_{\frac{\pi}{n+1}}^{\frac{\pi}{n}}


\varphi_c(x) dx \bigg)^{\frac{1}{p}};\\


%


\bigg(\int_{-\pi}^{\pi}\varphi(x)\bigg|\sum\limits_{n=1}^{\infty}


b_n\sin(2n-1)\frac{x}{2} \bigg|^p dx \bigg)^{\frac{1}{p}}\leq


C_2\bigg(\sum\limits_{n=1}^{\infty}b_n^p n^p 


\int_{\frac{\pi}{n+1}}^{\frac{\pi}{n}}


\varphi_c(x) dx \bigg)^{\frac{1}{p}}; \\


%


\bigg(\int_{-\pi}^{\pi}\varphi(x)\bigg|\sum\limits_{n=1}^{\infty}


b_n \sin nx\bigg|^p dx \bigg)^{\frac{1}{p}} \leq


C_3\bigg(\sum\limits_{n=1}^{\infty}


b_n^p n^p\int_{\frac{\pi}{n+1}}^{\frac{\pi}{n}}


\varphi_c(x)dx \bigg)^{\frac{1}{p}}; \\


%


\bigg( \int_{-\pi}^{\pi} \varphi (x) 


\bigg|\frac{b_0}{2}+\sum\limits_{n=1}^{\infty}


b_n \cos nx\bigg|^pdx \bigg)^{\frac{1}{p}} \leq


C_4\bigg(\int_{\frac{\pi}{2}}^{\pi}\varphi_c(x) dx b_0^p +


\sum\limits_{n=1}^{\infty} b_n^p n^p 


\int_{\frac{\pi}{n+1}}^{\frac{\pi}{n}}


\varphi_c(x) dx \bigg)^{\frac{1}{p}}; \\


%


\bigg( \int_{- \frac{\pi}{2}}^{\frac{\pi}{2}} 


\varphi(x)\bigg|\sum\limits_{n=1}^{\infty}


\Delta _1b_n \frac{\sin nx}{\sin x}


\bigg|^p dx \bigg)^{\frac{1}{p}} \leq


C_5\bigg( \int_{\frac{\pi}{4}}^{\frac{\pi}{2}}


\varphi_c(x) dx b_1^p + \sum\limits_{n=2}^{\infty} b_n^p n^p


\int_{\frac{\pi}{n+1}}^{\frac{\pi}{n}}


\varphi_c(x) dx \bigg)^{\frac{1}{p}},


\end{gather}


где положительные постоянные $C_1,\dotsc,C_5$


не зависят от $\{b_n\}$.


\end{lem}





\begin{pf}


Применяя преобразование Абеля (см., напр., аналогичные доказательства


в [1], cc.\,100, 651), получаем равенство


\begin{equation}


F(x) =\sum\limits_{n=0}^{\infty} \Delta _1 b_n B_n(x),


\end{equation}


где если


\begin{equation}


F(x)=\sum\limits_{n=1}^{\infty}


b_n \frac{\sin(2n-1)\frac{x}{2}}


{2\sin{\frac{x}{2}}},


\end{equation}


то


$ B_n(x)=\frac{\sin^2n{\frac{x}{2}}}


{\sin^2{\frac{x}{2}}}$, $\Delta _1b_0=0$;





если $F(x)=\sum\limits_{n=1}^{\infty}


b_n \sin(2n-1)\frac{x}{2}$,


то $ B_n(x)=\frac{\sin^2n{\frac{x}{2}}}{\sin\frac{x}{2}}$,


$\Delta _1 b_0=0$;





если


$F(x)=\sum\limits_{n=1}^{\infty} b_n \sin nx$,


то $ B_n(x)= \overline{B}{}_n^1(x)$,


$\Delta _1 b_0=0$, $B_0(x)=0$;





если


$F(x)=\frac{b_0}{2}+\sum\limits_{n=1}^{\infty} b_n\cos nx$,


то $ B_n(x)=B_n^1(x)$;





если


$F(x)=\sum\limits_{n=1}^{\infty} \Delta _1 b_n \frac{\sin nx}{\sin x}$,


то $ B_n(x)=\frac{\sin nx}{\sin x}$,


$\Delta _1 b_0=0$, $B_0(x)=0$.





Доказательство проведем лишь для случая (22). Рассмотрим


$$


J=\| F(x)\|^p_{((-\pi,\pi);p;\varphi )} 


= \int_{0}^{\pi} \varphi_c(x)|F(x)|^p dx.


$$


Применяя лемму 3, получаем


\begin{multline*}


J\leq C_1\bigg( \sum\limits_{m=0}^{\infty}


\int_{\frac{\pi}{2^{m+1}}}^{\frac{\pi}{2^m}}


\varphi_c(x)


\bigg|\sum\limits_{n=1}^{2^{m+1}-1} \Delta _1 b_n B_n(x)\bigg|^p dx+\\


%


+\sum\limits_{m=0}^{\infty}


\int_{\frac{\pi}{2^{m+1}}}^{\frac{\pi}{2^m}}


\varphi_c(x) \bigg|\sum\limits_{n=2^{m+1}}^{\infty} 


\Delta _1 b_n B_n(x)\bigg|^p dx\bigg) = C_1(J_1+J_2).


\end{multline*}





Так как $|B_n(x)|\leq C_2 n^2$, где положительная постоянная


$C_2$ не зависит от $x$ и $n$, то


$$


J_1\leq C_3 \sum\limits_{m=0}^{\infty} 


\int_{\frac{\pi}{2^{m+1}}}^{\frac{\pi}{2^m}}


\varphi_c(x) dx\bigg( \sum\limits_{n=1}^{2^{m+1}-1} 


\Delta _1 b_n(n+1)^2\bigg)^p.


$$


Поскольку $\Delta _1 b_n = b_n-b_{n+1}$ и $b_n\geq 0$, то


$$


\sum\limits_{n=1}^{2^m-1} \Delta _1 b_n(n+1)^2=\sum\limits_{n=2}^{2^m-1}


b_n[(n+1)^2-n^2]+b_1\cdot 2^2-b_{2^m}\cdot 2^m\leq 4


\sum\limits_{n=1}^{2^m-1} b_n\cdot n.


$$


Применяя эту оценку и учитывая свойства $\{b_n\}$, имеем


\begin{multline*}


J_1\leq C_4 \sum\limits_{m=0}^\infty{} 


\int_{\frac{\pi}{2^{m+1}}}^{\frac{\pi}{2^m}}


\varphi_c(x) dx\bigg( \sum\limits_{s=1}^{m+1} 


\sum\limits_{n=2^{s-1}}^{2^s-1} b_n\cdot n\bigg)^p\le\\


%


\leq C_5\bigg( \int_{\frac{\pi}{2}}^{\pi}


\varphi_c(x) dx\,b_1^p+\int_{\frac{\pi}{4}}^{\frac{\pi}{2}}


\varphi_c(x) dx\,b_1^p+


\sum\limits_{m=2}^{\infty} \int_{\frac{\pi}{2^m}}^{\frac{\pi}{2^{m-1}}}


\varphi_c(x) dx \bigg(\sum\limits_{s=2}^{m} 


b_{2^{s-1}}2^{2s}\bigg)^p\bigg).


\end{multline*}


Используя лемму 1 при $p\in[1,\infty)$ и лемму 2 при


$p\in(0,1)$, получаем


$$


J_1\leq C_6 \bigg( \int_{\frac{\pi}{2}}^{\pi} \varphi_c(x)


dx\,b_1^p+\sum\limits_{s=2}^{\infty} 


\int_{\frac{\pi}{2^{s-1}}}^{\frac{\pi}{2^{s-2}}}


\varphi_c(x) dx\,b_{2^{s-1}}^p 2^{sp}\bigg)\leq


C_7 \sum\limits_{n=1}^{\infty} b_n n^{2p}


\int_{\frac{\pi}{n+1}}^{\frac{\pi}{n}} \varphi_c(x) dx.


$$





Оценим теперь $J_2$. Так как $| B_n(x)|\leq C_8 / x^2$,


где положительная постоянная $C_8$ не зависит от $x$ и $n$, то


$$


J_2\leq C_8 \sum\limits_{m=0}^{\infty} 


\int_{\frac{\pi}{2^{m+1}}}^{\frac{\pi}{2^m}}


\varphi_c(x) dx 2^{2mp} \bigg| 


\sum\limits_{n=2^{m+1}}^{\infty} \Delta _1 b_m\bigg|^p\leq


C_9 \sum\limits_{n=1}^{\infty}


b_n n^{2p} \int_{\frac{\pi}{n+1}}^{\frac{\pi}{n}}


\varphi_c(x) dx.


$$


Объединяя оценки для $J_1$ и $J_2$, будем


иметь неравенство (16). Аналогично доказываются неравенства


(17)--(20).


\end{pf}





\begin{lem}


Пусть $\varphi \in\Phi$, а последовательность


$\{b_n\}$ является $(1,0)$-монотонной.


Тогда для $p\in(0,\infty)$





{\em а)} если $k=2$, то


\begin{align}


\bigg( \int_{-\pi}^{\pi}\varphi(x)\bigg|\sum\limits_{n=1}^{\infty} b_n


\frac{\sin(2n-1)\frac{x}{2}}{\sin\frac{x}{2}}


\bigg|^p dx \bigg)^{\frac{1}{p}} 


&\geq C_1\bigg( \sum\limits_{n=1}^{\infty} b_n^p n^{2p} 


\int_{\frac{\pi}{n+1}}^{\frac{\pi}{n}}


\varphi_c(x) dx\bigg)^{\frac{1}{p}}; \\


%


\bigg( \int_{-\frac{\pi}{2}}^{\frac{\pi}{2}} 


\varphi(x)\bigg|\sum\limits_{n=1}^{\infty} b_n 


\sin(2n-1)x \bigg|^p dx \bigg)^{\frac{1}{p}} 


&\geq C_2\bigg( \int_{\frac{\pi}{4}}^{\frac{\pi}{2}}


\varphi_c (x)dx\, b_1^p+S_2\bigg)^


{\frac{1}{p}}; \\


%


\bigg( \int_{-\pi}^{\pi}\varphi(x)\bigg|\sum\limits_{n=1}^{\infty} b_n


\sin nx \bigg|^p dx \bigg)^{\frac{1}{p}} &\geq


C_3\bigg( \int_{\frac{\pi}{2}}^{\pi} 


\varphi_c(x) dx\, b_1^p+S_2\bigg)^


{\frac{1}{p}};\\


%


\bigg( \int_{-\frac{\pi}{2}}^{\frac{\pi}{2}}


\varphi(x)\bigg|\sum\limits_{n=1}^{\infty} b_n


\frac{\sin nx}{\sin x} \bigg|^p dx \bigg)^{\frac{1}{p}} &\geq


C_4\bigg( \int_{\frac{\pi}{4}}^{\frac{\pi}{2}}


\varphi_c(x)dx\; b_1^p+S_2\bigg)^


{\frac{1}{p}};


\end{align}





{\em б)} если $k=1$ и, кроме того, функция $\varphi(x)$ дополнительно


удовлетворяет условию $A_2$, то справедливы неравенства $(23)$, $(24)$,


\begin{equation}


\bigg( \int_{-\pi}^{\pi}\varphi(x)\bigg|\frac{b_0}{2} + 


\sum\limits_{n=1}^{\infty} b_n


\sin nx \bigg|^p dx \bigg)^{\frac{1}{p}} 


\geq C_5\bigg( \int_{\frac{\pi}{4}}^{\frac{\pi}{2}}


\varphi_c (x)dx (b_0^p+b_1^p)+ S_2\bigg)^{\frac{1}{p}},


\end{equation}


где $S_2 = \sum\limits_{n=2}^{\infty} b_n^p n^p 


\dint_{\frac{\pi}{n+1}}^{\frac{\pi}{n}}\varphi_c (x) dx$,


а положительные постоянные $C_1,\dotsc, C_5$ не зависят от $\{b_n\}$.


\end{lem}





\begin{pf}


а) Докажем неравенство (23).


Используя представление (21), будем иметь 


$$


J = \|F(x)\|_{((-\pi,\pi);p;\varphi)}^{p} \geq C_1 \sum\limits_{v=1}^{\infty}


\int_{\frac{\pi}{v+1}}^{\frac{\pi}{v}}\varphi_c (x) \bigg|


\sum\limits_{n=v+1}^{\infty}\Delta_2 b_{n+1} B_n^3 (x) \bigg|^p dx.


$$


Как показано в [6], для $x\in [0,\pi]$ и


$n\geq\nu-1$ $B_n^3(x)\geq C_2\frac{n+1}{x^2}$. Поэтому 


$$


J\geq C_3 \sum\limits_{\nu =1}^{\infty}


\int_{\frac{\pi}{\nu +1}}^{\frac{\pi}{\nu}} \varphi_c(x) dx


\nu^{2p}


\bigg(\sum\limits_{n=\nu -1}^{\infty} \Delta _2 b_{n+1} (n+1)\bigg) ^p.


$$


Но $\sum\limits_{n= \nu}^{\infty} \Delta_2 b_n n =


\Delta_1 b_{\nu} \cdot \nu + b_{\nu +1} =


b_{\nu} + (\nu -1) \Delta_1 b_{\nu} \geq b_{\nu}$. 


Используя эту оценку, получим 


$J \geq C_3 \sum\limits_{\nu =1}^{\infty} b_{\nu}^p 


{\nu}^{2p} \dint_{\frac{\pi}{\nu +1}}^{\frac{\pi}{\nu}} 


\varphi_c (x) dx$. 


Неравенство (23) доказано. Аналогично доказываются неравенства


(24)--(26). Пункт~а) доказан.





б) Докажем неравенство (24) (случай $k=1$).


Используя представление (21), будем иметь\\


$$


\| F(x)\| ^p_{((-\frac{\pi}{2},\frac{\pi}{2});p;\varphi )}=


\sum\limits_{\nu=1}^{\infty}


\int_{\frac{\pi}{2^{\nu+1}}}^{\frac{\pi}{2^ \nu}}\varphi_c(x)


\bigg|\sum\limits_{n=1}^{\infty}


\Delta_1 b_n \frac{\sin^2 nx}{\sin x}\bigg|^p dx=


\sum\limits_{\nu=1}^{\infty} A_\nu,


$$


где $A_\nu = \dint_{\frac{\pi}{2^{\nu+1}}}^{\frac{\pi}{2^\nu}}


\varphi_c(x)|F(x)|^p dx$.


Покажем, что


\begin{equation}


A_\nu \geq C_4\cdot 2^{\nu p} b_{2^\nu -1}^p


\int_{\frac{\pi}{2^{\nu+1}}}^{\frac{\pi}{2^\nu}}


\varphi_c(x) dx,


\end{equation}


где положительная постоянная $C_4$


не зависит от $\nu$ и $\{b_\nu\}$.





Обозначим $J_\nu =[\frac{\pi}{2^{\nu+1}},\frac{\pi}{2^\nu}]$. Пусть


$J'_\nu = \{x\in J_\nu:\varphi_c(x)\neq 0\}$.


Если $\mu J'_\nu=0$, то неравенство(28) доказано. Пусть


$\mu J'_\nu>0$. Рассмотрим два случая.





1. Пусть


$\sum\limits_{k=1}^{2^\nu-1} k b_k >2\cdot 2^{2\nu }b_{2^\nu -1}$.


Для любого $x\in J_\nu$ имеем


\begin{multline*}


F(x) \geq \frac{4}{\pi^2} x \bigg( \sum\limits_{n=1}^{2^\nu-1}


b_n n^2- \sum\limits_{n=2}^{2^\nu} b_n (n-1)^2\bigg)\geq\\


%


\ge\frac{4}{\pi^2} x \bigg(


\sum\limits_{n=1}^{2^\nu-1} b_n\cdot n -\frac{1}{2} (2\cdot 2^{2\nu}b_{2^\nu


})\bigg)\geq


\frac{2}{\pi^2} \frac{\pi}{2^{\nu+1}} \sum\limits_{n=1}^{2^\nu-1}


b_n\cdot n\geq C_5\cdot 2^\nu b_{2^\nu -1}.


\end{multline*}


Тогда $A_\nu \geq C_6 2^{\nu p} b_{2^\nu-1}^{p}


\dint_{\frac{\pi}{2^{\nu+1}}}^{\frac{\pi}{2^\nu}}


\varphi_c(x) dx$,


где положительная постоянная $C_6$ не зависит от


$\nu $ и $\{b_{\nu} \}$.





2. Пусть $2\cdot 2^{2\nu } b_{2^\nu -1} \leq \sum\limits_{k=1}^{2^\nu-1}


k b_k$. Для любого $x\in J_\nu $ имеем


$$


F(x)= \sum\limits_{n=1}^{2^\nu -1} \Delta _1 b_n \frac{\sin^2 nx}{\sin x} +


\sum\limits_{n=2^\nu}^{\infty}


\Delta _1 b_n \frac{\sin^2 nx}{\sin x} \leq\\


\frac{\pi}{2} x \sum\limits_{n=1}^{2^\nu -1} 


\Delta _1 b_n\cdot n^2 +\frac{\pi}{2x}


\sum\limits_{n=2^\nu}^{\infty} \Delta _1 b_n \leq


C_7 2^\nu b_{2^\nu -1}.


$$





В то же время, используя условие $A_2$, получаем


\begin{multline*}


\int_{J'_\nu} \varphi_c(x) F(x) dx =


\int_{J_\nu} \varphi_c(x) F(x) dx \geq 


\sum\limits_{k=2^\nu-1}^{\infty} \Delta _1 b_k 


\int_{J_\nu} \frac{\sin^2 kx}{\sin x}


\varphi_c(x) dx \geq \\


%


\ge\sum\limits_{k=2^\nu-1}^{\infty} \Delta _1 b_k \frac{2^{\nu+1}}{\pi}


\int_{J_\nu} \sin^2 kx \varphi_c(x) dx\geq


C_8 2^\nu b_{2^\nu-1} \int_{J'_\nu} \varphi_c(x) dx.


\end{multline*}


Таким образом, \\


\begin{equation}


\int_{J'_\nu} \varphi_c(x) F(x) dx \geq


C_8 2^\nu b_{2^\nu-1} \int_{J'_\nu} \varphi_c(x) dx,


\end{equation}


где положительная постоянная $C_8$ не зависит от


$\nu$, $\{b_\nu\}$ и множества $J'_\nu$.





Пусть $I_\nu = \{x \in J'_\nu : F(x)\geq 


\frac{1}{2\pi} C_8 2^\nu b_{2^\nu-1}\}$.


Покажем, что


\begin{equation}


\int_{I_\nu} \varphi_c(x) dx \geq


\frac{C_8}{C_7} \frac{1}{4} \int_{J'_\nu} \varphi_c(x)dx.


\end{equation}


Докажем методом от противного. Предположим, что левая часть неравенства (30)


меньше правой. Тогда


\begin{multline*}


\int_{J'_\nu} \varphi_c(x) F(x) dx= \int_{I_\nu}


\varphi_c(x) F(x)dx + \int_{J'_\nu\setminus I_\nu} 


\varphi_c(x) F(x) dx \leq\\


%


\le C_7 2^\nu b_{2^\nu -1} \int_{I_\nu} \varphi_c(x) dx +


\frac{1}{2\pi} C_8 b_{2^\nu -1} \int_{J'_\nu} \varphi_c(x)


dx < \frac{1}{2} C_8 2^\nu b_{2^\nu - 1}


\int_{J'_\nu} \varphi_c(x) dx,


\end{multline*}


что противоречит неравенству (29). Следовательно, верно неравенство


(30). Но тогда


$$


A_\nu \geq \int_{I_\nu} 


\varphi_c(x) |F(x) |^p dx \geq C_9 2^{\nu p}


b_{2^\nu - 1}^p \int_{J_\nu} \varphi_c(x) dx.


$$


Таким образом, 


$$


\| F(x)\| ^p_{((-\frac{\pi}{2},\frac{\pi}{2});p;\varphi )}


\geq C_{10} \sum\limits_{\nu=1}^{\infty} 2^{\nu p} b_{2^\nu - 1}^p


\int_{\frac{\pi}{2^{\nu+1}}}^{\frac{\pi}{2^\nu}}\varphi _c(x) dx \geq 


C_{11} \sum\limits_{n=2}^{\infty} \int_{\frac{\pi}{n+1}}^{\frac{\pi}{n}}


\varphi_c(x) dx\,n^p b_{n-1}^p.


$$


Нетрудно проверить, что 


$$


\| F(x)\| ^p_{((-\frac{\pi}{2},


\frac{\pi}{2});p;\varphi )}


\geq C_{12} 2^p b_1^p \int_{\frac{\pi}{4}}^{\frac{\pi}{2}} 


\varphi_c(x) dx.


$$


Отсюда следует справедливость


неравенства (24). Аналогичным образом проверяются остальные неравенства.


Заметим только, что при доказательстве неравенства (27) используется


представление (15) функции $g(x)$ из леммы 4.


\end{pf}





\begin{lem}


Пусть последовательность $\{a_n\}$ такова, что


$a_n \to 0$ при $n \to \infty$ и


$ \Delta_k a_n \geq 0 $ для некоторого $k \geq 1$ и любого $n$.


Тогда для любого $l=0,1, \dotsc, k-1$ и любого


$n$ \ $\Delta_l a_n \geq 0 $.


\end{lem}





{\bf Доказательство } следует из свойств последовательности $\{a_n \}$


и определения $\Delta_k a_n$.$\quad\square$


\medskip





{\bf 3. Доказательство утверждения 2} 


состоит в последовательном применении


лемм 7, 4 и 5. Покажем это на примере п.\,a), случая $k\geq 2$, $s=0$.





Применяя лемму 7 и представление функции $f(x)$ для случая


$k=1$, $s=0$ из леммы 4, будем иметь 


\begin{equation}


f(x)=\frac{1}{2\sin {\frac{x}{2}}}


\sum\limits_{n=1}^{\infty} (a_{n-1} - a_n)


\sin {(2n-1) \frac{x}{2}} = \sum\limits_{n=1}^{\infty}


b_n \frac {\sin {(2n-1)\frac{x}{2}} }


{\sin \frac{x}{2}},


\end{equation}


где последовательность $\{b_n\}$ 


будет $(1,0)$-монотонной. Тогда, применяя


неравенство (16) из леммы 5, 


получаем справедливость неравенства (5). Аналогично


доказываются неравенства (6)--(8).$\quad\square$


\medskip





{\bf 4. Доказательство утверждения 3} 


состоит в последовательном применении лемм


7, 4, 6, 1 и 2. Покажем это на примере п.\,а), случая


$k \geq 3$, $s=0$. Применяя 


представление функции $f(x)$ из леммы 4, случай


$k=1$, $s=0$, имеем (31),


где последовательность $\{b_n\}$ будет $(2,0)$-монотонной.


Применяя неравенство $(23)$,


получаем справедливость (9). Пусть функция $\varphi (x)$


дополнительно удовлетворяет условию $A(p)$. Используя лемму 1 при


$p \in [1, \infty )$ и лемму 2 при $p \in (0,1)$, имеем


\begin{multline*}


\int_{\frac{\pi}{2}}^{\pi} \varphi_c(x)


dx\,a_0^p + \sum\limits_{n=1}^{\infty} a_n^p n^p


\int_{\frac{\pi}{n+1}}^{\frac{\pi}{n}}


\varphi_c(x) dx \le \\


%


\leq C_1 \sum\limits_{m=0}^{\infty} a_{[2^{m-1}]}^p 2^{(m-1)p}


\int_{\frac{\pi}{2^{m+1}}}^{\frac{\pi}{2^m}}


\varphi_c(x) dx \le \\


%


\leq C_2 \sum\limits_{m=0}^{\infty} \int_{\frac{\pi}


{2^{m+1}}}^{\frac{\pi}{2^m}}


\varphi_c(x) x^{-p} dx


\bigg( \sum\limits_{n=m}^{\infty} (a_{[2^{n-1}]} - a_{2^n})\bigg)^p \leq\\


%


\le C_3 \sum\limits_{m=0}^{\infty}


( a_{[2^{m-1}]} - a_{2^m})^p \int_{\frac{\pi}


{2^m}}^{\frac{\pi}{2^{m-1}}}


\varphi_c(x) x^{-p} dx \leq\\


%


\le C_4 \sum\limits_{n=1}^{\infty} ( \Delta_1 a_{n-1})^p n^{2p} 


\int_{\frac{\pi}{n+1}}^{\frac{\pi}{n}}


\varphi_c(x) dx \leq


C_5 \| f(x)\| ^p_{((-\pi,\pi);p;\varphi )}.


\end{multline*}


Таким образом, неравенство (11) доказано.


Аналогично доказываются неравенства (10), (12)--(14).$\quad\square$


\medskip





{\bf 5. Доказательство утверждения 4} основано на леммах 7, 4, 3 и 6. Покажем


это на примере п.\,а), случая


$k \geq 2$, $s=1$. Применяя лемму 7 и представление


функции $f(x)$ для случая $k=2$, $s=1$ из леммы 4, будем иметь


\begin{multline*}


f(x)= \frac {a_0 - a_2}{2} -


\frac {\Delta_1(\{\Delta\}_1 a_1) + 4 \Delta_1(\{\Delta\}_1 a_2)}{2}


\frac {1}{(2 \cos \frac{x}{2})^2} + \\


+2 \Delta_1(\{\Delta\}_1 a_2) \tg^2 \frac{x}{2} +


\frac{1}{(2 \cos \frac{x}{2})^2}


\sum\limits_{n=1}^{\infty} b_n\frac{\sin {(2n-1) 


\frac{x}{2}}}{2 \sin \frac{x}{2}},


\end{multline*}


где $b_1=b_2= \Delta_1(\{\Delta\}_1 a_1) + 


\Delta_1(\{\Delta\}_1 a_2)$,


$b_n= \Delta_1(\{\Delta\}_1 a_{n-2})$, $n=3,4,\dotsc$,


последовательность $\{b_n\}$ является


$(1,0)$-монотонной.


При $x \in [ \frac{\pi}{4}, \frac{\pi}{2} ]$ 


и условии $a_0 \geq a_2 + a_1 - a_3 + 4(a_2 - a_4)$, 


получаем 


\begin{multline*}


\frac{a_0-a_2}{2} - 


\frac{\Delta_1(\{\Delta \}_1 a_1)+4 \Delta_1(\{ \Delta \}_1 a_2)}{2}


\frac{1}{(2 \cos {\frac{x}{2}})^2} +


\frac{\Delta_1(\{ \Delta \}_1 a_2)}{2} \tg^2 {\frac{x}{2}} \geq \\


%


\ge\frac{a_0 - a_2}{2} -


\frac{\Delta_1 (\{ \Delta \}_1 a_1) + 4\Delta_1 (\{ \Delta \}_1 a_2)}{4}


+ \frac{\Delta_1 (\{ \Delta \}_1 a_2)}{12} \geq \\


%


\ge \frac{2(a_0 - a_2) - (\Delta_1 ( \{ \Delta \}_1 a_1) +


4 \Delta_1 (\{ \Delta \}_1 a_2))}{4} \geq


\frac{\Delta_1 (\{ \Delta \}_1 a_0) }{4}.


\end{multline*}


Учитывая последнюю оценку


и применяя леммы 3 и 6, окончательно получим


(10). Неравенство (12) в случае, 


когда функция $ \varphi (x)$


дополнительно удовлетворяет условию $A(p)$, 


доказывается так же, как проводилось доказательство


утверждения 3 в этом случае. 


Аналогично доказываются остальные неравенства


пп. а) и б).$\quad\square$





\begin{note}


Утверждения 2, 3, 4  сформулированы для $s=0, 1, 2$.


Аналогично рассуждая, можно получить подобные утверждения и для


случая, когда последовательность $\{a_n \} $ будет $(k,s)$-монотонной


при $s>2$.


\end{note}





\begin{note}


Пусть $ \varphi (x) \equiv 1$, а в (3) $a_0=0$. Тогда эта 


функция удовлетворяет условиям, наложенным на функцию $ \varphi (x)$


в утверждении 2. Из (4) следует, что $ \varphi (x) \equiv 1$


удовлетворяет условию $A_2$ в утверждении 4. В силу этого


утверждения 2 и 4 являются обобщениями теоремы из работы [2].


\end{note}





\begin{note}


При более жестких ограничениях на функцию $\varphi (x)$


результаты утверждения~2 можно перенести с интервала


$(- \frac {\pi}{2}, \frac{\pi}{2})$ на $(- \pi,\pi)$.





Так, например,


пусть $\varphi$ из $\Phi$, четна и удовлетворяет условиям $A_1$,


$A_2$, $A_3$; последовательность $\{a_n\}$ является $(k,s)$-монотонной.


\end{note}





Представим функции $f(x)$ и $g(x)$ в виде


\begin{align*}


f(x+ \pi) &= f(x) - 2 \sum\limits_{k=1}^{\infty} a_{2k-1} \cos (2k-1)x,\\


g(x+ \pi) &= g(x) - 2 \sum\limits_{k=1}^{\infty} a_{2k-1} \sin (2k-1)x.


\end{align*}


Тогда


\begin{align*}


\int_{-\pi}^{\pi} \varphi (x) | f(x) |^p dx &\leq


C_1 \bigg( \int_{-\frac{\pi}{2}}^{\frac{\pi}{2}} 


\varphi (x) | f(x)|^p dx + \\


&\qquad+\int_{-\frac{\pi}{2}}^{\frac{\pi}{2}} 


\varphi (x) | \sum\limits_{k=1}^{\infty}


a_{2k-1} \cos (2k-1)x |^p dx \bigg),\\


%


\int_{-\pi}^{\pi} \varphi (x) | g(x) |^p dx &\leq


C_2 \bigg( \int_{-\frac{\pi}{2}}^{\frac{\pi}{2}} 


\varphi (x) | g(x)|^p dx + \\


&\qquad +\int_{-\frac{\pi}{2}}^{\frac{\pi}{2}} 


\varphi (x) \bigg| \sum\limits_{k=1}^{\infty}


a_{2k-1} \sin (2k-1)x \bigg|^p dx \bigg).


\end{align*}





Пусть последовательность $\{a_n\}$ \ $(1,1)$-монотонна. Тогда для


$p \in (0, \infty)$ из утверждений 2 и 4 имеем 


\begin{gather*}


\| g(x)\| _{((-\pi,\pi);p;\varphi )}


\leq C_3 \bigg( \int_{\frac{\pi}{4}}^{\frac{\pi}{2}} \varphi_c(x)


dx (\{\Delta\}_1 a_0)^p + 


\sum\limits_{n=2}^{\infty} (\{\Delta\}_1 a_{n-1})^p n^p


\int_{\frac{\pi}{n+1}}^{\frac{\pi}{n}}


\varphi_c(x) dx\bigg)^{\frac{1}{p}} = C_3 \cdot B,\\


\| g(x)\| _{((-\pi,\pi);p;\varphi )} \geq C_4\cdot B,\quad


\| f(x)\| _{((-\pi,\pi);p;\varphi )} \leq C_5\cdot B.


\end{gather*}





Пусть $f(x)= \frac{a_1 - a_3}{4} + \frac{a_0}{2} +


\sum\limits_{n=1}^{\infty} a_n \cos nx$,


где $\Delta_2(\{\Delta\}_2 a_n) \geq 0$, $n= 0, 1,\dotsc$


Тогда, применяя утверждения 2, 4 и используя представление


$$


f(x)= \frac{a_1 - a_3}{4} - \frac{a_1 - a_3}{2}


\frac{1}{(2 \cos \frac{x}{2})^2} + \frac{1}{(2 \cos \frac{x}{2})^2}


\bigg(\frac{2a_0+3a_1-a_3}{2} + \\


\sum\limits_{n=1}^{\infty} \{ \Delta \}_2 a_{n-1} \cos {nx}\bigg),


$$


имеем


\begin{gather*}


\| f(x)\| _{((-\pi,\pi);p;\varphi )} \leq


C_5 \bigg( \int_{\frac{\pi}{4}}^{\frac{\pi}{2}} \varphi_c(x)


dx ((\Delta_1(\{\Delta\}_2 a_0))^p +(\Delta_1(\{\Delta\}_2 a_1))^p) +\\


+\sum\limits_{n=4}^{\infty} \int_{\frac{\pi}{n+1}}^{\frac{\pi}{n}}


\varphi_c(x) dx\,n^{2p} (\Delta_1(\{\Delta\}_2 a_{n-2})


)^p\bigg)^{\frac{1}{p}} = C_5\cdot B_1,\\


\| f(x)\|_{((-\pi,\pi);p;\varphi )} \geq C_6 \cdot B_1.


\end{gather*}





{\bf 7. Доказательство утверждения 1 }


следует из (1), утверждений 2, 4 и замечания 3,


если взять функцию $\varphi (x) \equiv 1$ 


(она удовлетворяет всем требуемым


условиям).$\quad\square$








\begin{thebibliography}{9}





\bibitem{1}


Бари Н.К. {\it Тригонометрические ряды.} --


М.: Физматгиз, 1961. -- 936 с. 


\bibitem{2}


Вуколова Т.М., Дьяченко М.И. {\it 


О свойствах сумм тригонометрических рядов с


монотонными коэффициентами} // Вестн. Моск. ун-та. Сер. матем., механ. --


1995. -- \No\,3. -- С.\,22--32.


\bibitem{3}


Бари Н.К., Стечкин С.Б. {\it Наилучшие приближения и дифференциальные


свойства двух сопряженных функций} // Тр. Моск. матем. о-ва. --


1956. -- Т.\,5. -- С.\,483--522. 


\bibitem{4}


Потапов М.К., Бериша М. {\it Модули гладкости и коэффициенты Фурье


периодических функций одного переменного} // Publ. Inst. math. --


1979. -- T.\,26. -- P.\,215--228. 


\bibitem{5}


Никольский С.М. {\it Приближение функций многих переменных и теоремы


вложения}. -- М.: Наука, 1977. -- 455 с. 


\bibitem{6}


Вуколова Т.М. {\it О рядах по синусам и косинусам с кратно


монотонными коэффициентами} // Вестн. Моск. ун-та. Сер. матем.,


механ. -- 1990. -- \No\,5. -- С.\,38--42. 





\end{thebibliography}





\vskip 2cm





\postup{\it Волгоградский государственный}{\qquad\it Поступили}


\postup{\it технический университет}{\it первый вариант $29.05.2001$}


\postup{}{\it окончательный вариант $19.08.2002$}


\label{end}


\end{document}





